\chapter{Connections and Curvature}\label{connections}
\chapterstatus{complete}

\section{Introduction}\label{connections:section-introduction}

To differentiate a function $f$ on a manifold along a curve $\gamma$ one compares the values $f(\gamma(t))$ and $f(\gamma(0))$, which are numbers. For a section $s$ of a vector
bundle $E$ the values $s(\gamma(t)) \in E_{\gamma(t)}$ and $s(\gamma(0)) \in E_{\gamma(0)}$ lie in \emph{different} vector spaces, and there is no canonical way to subtract
them. A local trivialisation identifies nearby fibres with $\mathbb{K}^r$, and then one can differentiate the components of $s$; but the result depends on the trivialisation,
because changing it multiplies the components by a matrix of functions, whose derivative enters the answer. The simplest example is a vector field $Y$ on the sphere
$S^2 \subseteq \RR^3$: its derivative in $\RR^3$ along a curve on the sphere is in general not tangent to the sphere, so it is not a candidate for a derivative of $Y$ within
the tangent bundle of $S^2$.

A \emph{connection} is the extra structure that removes this ambiguity: a rule $\nabla$ that assigns to a section $s$ and a tangent vector $v$ a derivative $\nabla_vs$ in the
same fibre, subject to the Leibniz rule. Everything else follows from it. Solving $\nabla_{\dot\gamma}s = 0$ along curves gives \emph{parallel transport}, a way of moving
vectors from fibre to fibre along a path, which in general depends on the path; going around loops gives the \emph{holonomy group}. The \emph{curvature} measures the failure
of covariant derivatives to commute, $F(X, Y) = \nabla_X\nabla_Y - \nabla_Y\nabla_X - \nabla_{[X, Y]}$; unlike the connection itself it is a tensor, an $\End(E)$-valued
$2$-form, and it is the infinitesimal form of the holonomy around small loops (Proposition~\ref{connections:proposition-small-loops}). Connections with zero curvature are the
same as local systems (Theorem~\ref{connections:theorem-flat-local-systems}), which is the differential-geometric form of the monodromy correspondence of
Chapter~\ref{sheaves} and the starting point of the theory of variations of Hodge structure (Chapter~\ref{vhs}), where the Gauss--Manin connection appears. The curvature is
also the source of the characteristic classes of Chapter~\ref{characteristic-classes}, and for the tangent bundle of a Riemannian manifold it is the Riemann curvature tensor
of Chapter~\ref{riemannian-geometry}.

The chapter goes as follows.
\begin{enumerate}
\item Forms with values in a vector bundle, the objects in which connections and curvature take their values; the tensoriality lemma, which says when an operator on
sections is given pointwise (Section~\ref{connections:section-bundle-valued-forms}).
\item Connections, in the two equivalent languages $\nabla \colon \Gamma(E) \to \Omega^1(M, E)$ and $(X, s) \mapsto \nabla_Xs$; connection forms in local frames and how they
change; existence; the basic examples, among them the connection of a subbundle of a trivial bundle; induced connections on duals, tensor products and
$\Hom$ bundles (Section~\ref{connections:section-connections}).
\item The exterior covariant derivative and the curvature, with complete computations for the sphere and for the tautological line bundle of $\CC\PP^1$
(Section~\ref{connections:section-curvature}).
\item Parallel transport, holonomy, the geometric meaning of curvature, and flat connections (Section~\ref{connections:section-parallel-transport}).
\item Metric connections and principal bundles (Sections~\ref{connections:section-metric} and~\ref{connections:section-principal}).
\end{enumerate}
References: \cite{kobayashi-nomizu}, \cite{milnor-stasheff}, \cite[Chapter 0]{griffiths-harris}.

\noindent\textbf{Prerequisites.} Chapter~\ref{vector-bundles} (Vector Bundles), Chapter~\ref{differential-forms} (Differential Forms
and Integration) and Chapter~\ref{lie-groups} (Lie Groups).

\noindent\textbf{Conventions.} $\mathbb{K}$ is $\RR$ or $\CC$, $M$ is a smooth manifold of dimension $n$, and $E \to M$ is a $\mathbb{K}$-vector bundle of rank $r$. For an open
$U \subseteq M$, $\Gamma(U, E)$ is the space of smooth sections over $U$, and $\Gamma(E) = \Gamma(M, E)$; $\mathfrak{X}(M)$ is the space of vector fields. A \emph{frame} of $E$
over $U$ is a tuple $e = (e_1, \ldots, e_r)$ of sections that is a basis of every fibre (Definition~\ref{vector-bundles:definition-sections}). When $\mathbb{K} = \CC$,
differential forms have complex values: $\Omega^k(M)$ then means $\Omega^k(M) \otimes_\RR \CC$, and $d$ is extended $\CC$-linearly. We write sums explicitly, and in
computations with frames we use matrix notation: a section $s = \sum_js^je_j$ is identified with the column vector $(s^1, \ldots, s^r)^T$ of its components.

\section{Bundle-valued forms}\label{connections:section-bundle-valued-forms}

A connection takes a section $s$ and produces, at each point $p$, a linear map $T_pM \to E_p$, $v \mapsto \nabla_vs$, the derivative of $s$ in the direction $v$; the curvature
takes two tangent vectors and produces an endomorphism of $E_p$. Such objects are forms with values in a vector bundle, and this section explains what they are, in several
equivalent ways.

\begin{definition}\label{connections:definition-bundle-valued-forms}
The space of \emph{$E$-valued $k$-forms} on $M$ is
\[
\Omega^k(M, E) = \Gamma(M, \Lambda^kT^*M \otimes E),
\]
the smooth sections of the tensor product of the bundle of $k$-forms with $E$ (Proposition~\ref{vector-bundles:proposition-operations}). So $\Omega^0(M, E) = \Gamma(M, E)$,
and for the trivial line bundle $\Omega^k(M, M \times \mathbb{K}) = \Omega^k(M)$. We write $\Omega^k(U, E)$ for the forms over an open set $U$.
\end{definition}

The following lemma is used constantly. It explains when an operator that takes sections to sections is really a pointwise, linear algebra operation.

\begin{lemma}[Tensoriality]\label{connections:lemma-tensoriality}
Let $E_1, \ldots, E_k$ and $F$ be vector bundles over $M$, and let $A \colon \Gamma(E_1) \times \cdots \times \Gamma(E_k) \to \Gamma(F)$ be $C^\infty(M)$-linear in each
argument: $A(\ldots, fs + gt, \ldots) = fA(\ldots, s, \ldots) + gA(\ldots, t, \ldots)$ for functions $f$, $g$. Then the value $A(s_1, \ldots, s_k)(p)$ depends only on the
vectors $s_1(p), \ldots, s_k(p)$. So there are multilinear maps $A_p \colon (E_1)_p \times \cdots \times (E_k)_p \to F_p$ with $A(s_1, \ldots, s_k)(p) = A_p(s_1(p), \ldots, s_k(p))$,
and they form a smooth section of $\Hom(E_1 \otimes \cdots \otimes E_k, F)$.
\end{lemma}

\begin{proof}
\emph{Locality.} Suppose $s_1$ vanishes on an open set $U \ni p$. By Lemma~\ref{smooth-manifolds:lemma-bump-manifold} there is a function $\chi$ with $\chi(p) = 1$ and support
in $U$. Then $\chi s_1 = 0$, so $0 = A(\chi s_1, s_2, \ldots) = \chi A(s_1, s_2, \ldots)$, and evaluating at $p$ gives $A(s_1, \ldots)(p) = 0$.

\emph{Pointwise dependence.} Suppose $s_1(p) = 0$. Let $(e_1, \ldots, e_m)$ be a frame of $E_1$ on a neighbourhood $U$ of $p$, and write $s_1 = \sum_ja^je_j$ on $U$ with smooth
$a^j$ vanishing at $p$. Choose $\chi$ as above, equal to $1$ on a neighbourhood $V$ of $p$, with support in $U$. The sections $\tilde e_j = \chi e_j$ and the functions
$\tilde a^j = \chi a^j$, extended by $0$ outside $U$, are smooth on $M$, and $\sum_j\tilde a^j\tilde e_j = \chi^2s_1$ agrees with $s_1$ on $V$. By locality and linearity,
$A(s_1, \ldots)(p) = A(\sum_j\tilde a^j\tilde e_j, \ldots)(p) = \sum_j\tilde a^j(p)A(\tilde e_j, \ldots)(p) = 0$. By linearity in each argument, $A(s_1, \ldots, s_k)(p)$ depends only
on the values $s_i(p)$.

\emph{The maps $A_p$.} Every vector $v \in (E_i)_p$ is the value at $p$ of a global section, for instance of $\chi\sum_jv^je_j$ with constant coefficients $v^j$. So
$A_p(v_1, \ldots, v_k) = A(s_1, \ldots, s_k)(p)$, for any sections with $s_i(p) = v_i$, is well defined and multilinear. In local frames of the $E_i$ and of $F$, the matrix
coefficients of $A_p$ are the components of $A$ applied to (cut-off) frame sections, which are smooth; so $p \mapsto A_p$ is a smooth section of
$\Hom(E_1 \otimes \cdots \otimes E_k, F)$.
\end{proof}

\begin{proposition}\label{connections:proposition-bundle-valued-forms}
For an $E$-valued $k$-form $\omega$ and a point $p$, the value $\omega_p \in \Lambda^kT_p^*M \otimes E_p$ is the same as an alternating $k$-linear map $(T_pM)^k \to E_p$, by
$(\alpha \otimes e)(v_1, \ldots, v_k) = \alpha(v_1, \ldots, v_k)\,e$. Moreover the following descriptions of $\Omega^k(M, E)$ are equivalent.
\begin{enumerate}
\item Smooth sections $\omega$ of $\Lambda^kT^*M \otimes E$.
\item Families of alternating $k$-linear maps $\omega_p \colon (T_pM)^k \to E_p$ such that, for all vector fields $X_1, \ldots, X_k$ on an open set $U$, the section
$p \mapsto \omega_p(X_1(p), \ldots, X_k(p))$ of $E$ over $U$ is smooth.
\item Maps $\omega \colon \mathfrak{X}(M)^k \to \Gamma(E)$ that are alternating and $C^\infty(M)$-linear in each argument.
\item Locally, for a frame $(e_1, \ldots, e_r)$ of $E$ over $U$: expressions $\omega|_U = \sum_i\omega^i \otimes e_i$ with ordinary $k$-forms $\omega^i \in \Omega^k(U)$, which
are uniquely determined by $\omega$ and are smooth exactly when $\omega$ is. Equivalently $\omega(X_1, \ldots, X_k) = \sum_i\omega^i(X_1, \ldots, X_k)\,e_i$.
\end{enumerate}
\end{proposition}

\begin{proof}
\emph{Pointwise.} For finite-dimensional $V$ and $W$, $\Lambda^kV^* \otimes W \cong \Hom(\Lambda^kV, W)$ by Propositions~\ref{multilinear-algebra:proposition-dual-exterior}
and~\ref{multilinear-algebra:proposition-hom-tensor}, and linear maps $\Lambda^kV \to W$ are the same as alternating $k$-linear maps $V^k \to W$ by the universal property of
the exterior power (Proposition~\ref{multilinear-algebra:proposition-universal-algebras}). With the convention of Notation~\ref{multilinear-algebra:notation-forms}, $\alpha \otimes e$
corresponds to $(v_1, \ldots, v_k) \mapsto \alpha(v_1, \ldots, v_k)e$.

(1) $\Leftrightarrow$ (4). If $(e_i)$ is a frame of $E$ over $U$ and $(dx^I)$ the frame of $\Lambda^kT^*M$ given by a chart on $U$, the products $dx^I \otimes e_i$ form a frame
of $\Lambda^kT^*M \otimes E$ over $U$ (Proposition~\ref{vector-bundles:proposition-operations}), and a section is smooth if and only if its coefficients $\omega^i_I$ in this frame are smooth.
Grouping, $\omega = \sum_i\omega^i \otimes e_i$ with $\omega^i = \sum_I\omega^i_I\,dx^I$; the $\omega^i$ are unique because at each point
$\Lambda^kT_p^*M \otimes E_p = \bigoplus_i\Lambda^kT_p^*M \otimes \mathbb{K}e_i(p)$, and $\omega$ is smooth exactly when all the $\omega^i$ are.

(2) $\Leftrightarrow$ (4). In a frame, $\omega(X_1, \ldots, X_k) = \sum_i\omega^i(X_1, \ldots, X_k)e_i$, which is smooth for all vector fields if and only if every
$\omega^i(X_1, \ldots, X_k)$ is; by Proposition~\ref{differential-forms:proposition-forms-descriptions} this holds if and only if the $\omega^i$ are smooth forms.

(1) $\Leftrightarrow$ (3). A form $\omega$ gives $(X_1, \ldots, X_k) \mapsto \omega(X_1, \ldots, X_k)$, which is alternating and $C^\infty(M)$-linear because it is computed
pointwise. Conversely, a map as in (3) is given pointwise by alternating multilinear maps $\omega_p$ (Lemma~\ref{connections:lemma-tensoriality}, applied with $E_i = TM$),
which satisfy (2).
\end{proof}

\begin{example}\label{connections:example-bundle-valued-forms}
\begin{enumerate}
\item On the trivial bundle $M \times \mathbb{K}^r$ with its standard frame, an $E$-valued $k$-form is a column vector $(\omega^1, \ldots, \omega^r)^T$ of ordinary $k$-forms.
\item A $TM$-valued $1$-form is at each point a linear map $T_pM \to T_pM$: $\Omega^1(M, TM)$ is the space of endomorphisms of the tangent bundle. The identity is the
$TM$-valued $1$-form $\sum_idx^i \otimes \partial/\partial x^i$ in every chart; the shape operator of a hypersurface (Chapter~\ref{riemannian-geometry}) is another example.
\item An $\End(E)$-valued $k$-form takes $k$ tangent vectors to an endomorphism of the fibre. In a frame $(e_i)$ it is an $r \times r$ matrix $(\omega^i_j)$ of $k$-forms, with
$\omega(X_1, \ldots, X_k)e_j = \sum_i\omega^i_j(X_1, \ldots, X_k)e_i$. The curvature of a connection will be such a form with $k = 2$.
\item The differential of a smooth map $f \colon N \to M$ is at each point a linear map $T_qN \to T_{f(q)}M = (f^*TM)_q$, so $df$ is an $f^*TM$-valued $1$-form on $N$.
\end{enumerate}
\end{example}

\emph{Products.} Forms with values in bundles can be multiplied whenever the values can: the wedge product of an ordinary $k$-form and an $E$-valued $l$-form is the $E$-valued
$(k + l)$-form given by $\alpha \wedge (\beta \otimes s) = (\alpha \wedge \beta) \otimes s$; an $\End(E)$-valued form acts on an $E$-valued form by
$(\alpha \otimes A) \wedge (\beta \otimes s) = (\alpha \wedge \beta) \otimes As$; and two $\End(E)$-valued forms multiply by
$(\alpha \otimes A) \wedge (\beta \otimes B) = (\alpha \wedge \beta) \otimes AB$. These formulas define bilinear maps fibre by fibre, hence operations on forms. In a frame they are
matrix products in which the entries are multiplied with the wedge product:
\[
(\theta \wedge \omega)^i = \sum_j\theta^i_j \wedge \omega^j, \qquad (\theta \wedge \eta)^i_j = \sum_k\theta^i_k \wedge \eta^k_j .
\]
We also write $\alpha \wedge s = \alpha \otimes s$ for $s \in \Omega^0(M, E)$. On vectors, by Proposition~\ref{multilinear-algebra:proposition-wedge-formula}, for $1$-forms
$\theta \in \Omega^1(M, \End E)$ and $\omega \in \Omega^1(M, E)$,
\[
(\theta \wedge \omega)(X, Y) = \theta(X)\omega(Y) - \theta(Y)\omega(X), \qquad (\theta \wedge \theta)(X, Y) = \theta(X)\theta(Y) - \theta(Y)\theta(X) = [\theta(X), \theta(Y)] .
\]
In particular, for a matrix $\theta$ of $1$-forms, $\theta \wedge \theta$ need not vanish: for $\theta = \bigl(\begin{smallmatrix} 0 & dx\\ dy & 0\end{smallmatrix}\bigr)$,
$\theta \wedge \theta = \bigl(\begin{smallmatrix} dx \wedge dy & 0\\ 0 & dy \wedge dx\end{smallmatrix}\bigr)$. For a single $1$-form ($r = 1$), $\theta \wedge \theta = 0$. The graded
commutativity of ordinary forms, $\alpha \wedge \beta = (-1)^{kl}\beta \wedge \alpha$, holds for an ordinary form $\alpha$ against any bundle-valued form $\beta$, but not for two
matrix-valued forms.

\begin{remark}\label{connections:remark-no-canonical-d}
There is no canonical exterior derivative on $\Omega^k(M, E)$. On a trivial bundle one can differentiate components, $d(\sum_i\omega^ie_i) = \sum_id\omega^i\,e_i$; but this
depends on the frame. If $e^{\prime} = eg$ is another frame, that is $e^{\prime}_j = \sum_ie_ig^i_j$ for a smooth $g \colon U \to \GL_r(\mathbb{K})$, then the components are
related by $\omega = g\omega^{\prime}$ (as columns: $\omega^i = \sum_jg^i_j\omega^{\prime j}$), and
\[
\sum_id\omega^i\,e_i = \sum_jd\omega^{\prime j}\,e^{\prime}_j + \sum_{i,j}(dg^i_j \wedge \omega^{\prime j})\,e_i .
\]
The two ``derivatives'' agree for all $\omega$ only if $dg = 0$, that is, if the change of frame is locally constant. A connection supplies the missing correction term, and it
does so consistently for all frames.
\end{remark}

\section{Connections}\label{connections:section-connections}

\begin{definition}\label{connections:definition-connection}
A \emph{connection} (or \emph{covariant derivative}) on a vector bundle $E \to M$ is a map $\nabla \colon \Gamma(M, E) \to \Omega^1(M, E)$ that is $\mathbb{K}$-linear and satisfies
the Leibniz rule
\[
\nabla(fs) = df \otimes s + f\nabla s, \qquad f \in C^\infty(M),\ s \in \Gamma(M, E).
\]
For a vector field $X$ we write $\nabla_Xs = (\nabla s)(X)$, so $\nabla_{fX}s = f\nabla_Xs$ and $\nabla_X(fs) = (Xf)s + f\nabla_Xs$.
\end{definition}

So $\nabla s$ is an $E$-valued $1$-form: at each point it is the linear map $T_pM \to E_p$, $v \mapsto \nabla_vs = (\nabla s)_p(v)$, the \emph{covariant derivative} of $s$ in the
direction $v$. In Riemannian geometry and in many books a connection is introduced instead as an operation $(X, s) \mapsto \nabla_Xs$, and for $E = TM$ one writes
$\nabla_XY$. The two points of view are equivalent.

\begin{proposition}\label{connections:proposition-connection-operators}
Let $\nabla$ be a connection. The map $\mathfrak{X}(M) \times \Gamma(E) \to \Gamma(E)$, $(X, s) \mapsto \nabla_Xs$, has the properties
\begin{enumerate}
\item $\nabla_{fX + gY}s = f\nabla_Xs + g\nabla_Ys$ for functions $f$, $g$: it is $C^\infty(M)$-linear in $X$;
\item $\nabla_X(s + t) = \nabla_Xs + \nabla_Xt$ and $\nabla_X(cs) = c\nabla_Xs$ for $c \in \mathbb{K}$;
\item $\nabla_X(fs) = (Xf)s + f\nabla_Xs$.
\end{enumerate}
Conversely, every map $(X, s) \mapsto D_Xs$ with these three properties is $D_Xs = \nabla_Xs$ for a unique connection $\nabla$. In particular $(\nabla_Xs)(p)$ depends only on
$X(p)$, so $\nabla_vs$ is defined for a single tangent vector $v$.
\end{proposition}

\begin{proof}
(1) holds because $\nabla s$ is a $1$-form, (2) is the linearity of $\nabla$, and (3) is the Leibniz rule evaluated on $X$, since $(df \otimes s)(X) = df(X)s = (Xf)s$.
Conversely, for fixed $s$ the map $X \mapsto D_Xs$ is $C^\infty(M)$-linear by (1), so by Proposition~\ref{connections:proposition-bundle-valued-forms}(3) it is an $E$-valued
$1$-form $\nabla s$ with $(\nabla s)(X) = D_Xs$. The map $s \mapsto \nabla s$ is linear by (2), and (3) says $\nabla(fs)(X) = (df \otimes s + f\nabla s)(X)$ for all $X$. The last
statement is Lemma~\ref{connections:lemma-tensoriality} for $X \mapsto \nabla_Xs$.
\end{proof}

Note the asymmetry: $\nabla_Xs$ is $C^\infty$-linear in $X$ but not in $s$. It depends only on the value of $X$ at a point, but on the values of $s$ near the point (in fact, on
its values along a curve with velocity $X(p)$; Remark~\ref{connections:remark-along-curves}).

\begin{example}\label{connections:example-connections}
\begin{enumerate}
\item \emph{The trivial connection.} On the trivial bundle $M \times \mathbb{K}^r$ a section is a vector-valued function $s = (s^1, \ldots, s^r)^T$, and
$ds = (ds^1, \ldots, ds^r)^T$ is a connection: $d(fs) = df\,s + f\,ds$. Here $\nabla_Xs = X(s)$ is the componentwise derivative.
\item \emph{Adding a matrix of $1$-forms.} For an $r \times r$ matrix $A$ of $1$-forms on $M$, $\nabla s = ds + As$ is a connection on $M \times \mathbb{K}^r$, since the extra term
is $C^\infty(M)$-linear in $s$: $A(fs) = fAs$. We will see that every connection on a trivial bundle has this form.
\item \emph{The Euclidean connection.} On $T\RR^n = \RR^n \times \RR^n$ the trivial connection is $\nabla_XY = DY(X) = \sum_iX(Y^i)\partial_i$, the directional derivative of
the components of $Y$.
\item \emph{Subbundles of trivial bundles.} Let $E \subseteq M \times \RR^N$ be a subbundle, so each $E_p$ is a subspace of $\RR^N$, and let $P_p \colon \RR^N \to E_p$ be the
orthogonal projection. Locally $E$ has a frame, and Gram--Schmidt turns it into an orthonormal frame $(u_1, \ldots, u_r)$; then $P = \sum_iu_iu_i^T$, so $P$ is a smooth
matrix-valued function. A section of $E$ is a function $s \colon M \to \RR^N$ with $s(p) \in E_p$, and
\[
\nabla s = P(ds), \qquad\text{that is,}\qquad \nabla_Xs = P(X(s)),
\]
is a connection: $\nabla(fs) = P(df\,s + f\,ds) = df \otimes s + fP(ds)$, since $Ps = s$. For the tangent bundle $TS^n \subseteq S^n \times \RR^{n+1}$ this says: differentiate
the vector field in $\RR^{n+1}$ and project back onto the tangent space. This is the Levi-Civita connection of the round sphere (Chapter~\ref{riemannian-geometry}). The same
construction applies to the tautological line bundles of $\RR\PP^n$ and $\CC\PP^n$ (with the Hermitian projection $P(w) = \sum_i\langle w, u_i\rangle u_i$ in the complex case),
to the tangent and normal bundles of any submanifold of $\RR^N$, and to the tautological bundles of Grassmannians.
\end{enumerate}
\end{example}

\begin{proposition}\label{connections:proposition-connections-exist}
Let $\nabla$ be a connection on $E$.
\begin{enumerate}
\item \emph{Locality.} If $s$ vanishes on an open set $U$, so does $\nabla s$. Consequently $\nabla$ induces a connection on $E|_U$ for every open $U$, with
$(\nabla s)|_U = \nabla(s|_U)$.
\item \emph{Connection forms.} If $e = (e_1, \ldots, e_r)$ is a frame over $U$, there is a unique matrix $\theta = (\theta^i_j)$ of $1$-forms on $U$, the \emph{connection form} of
$\nabla$ in the frame $e$, with $\nabla e_j = \sum_i\theta^i_j \otimes e_i$. For $s = \sum_js^je_j$,
\[
\nabla s = \sum_i\Bigl(ds^i + \sum_j\theta^i_js^j\Bigr) \otimes e_i, \qquad\text{in vector notation}\quad \nabla s = ds + \theta s .
\]
\item \emph{Change of frame.} If $e^{\prime} = eg$ is another frame over $U$, with $g \colon U \to \GL_r(\mathbb{K})$ smooth, then
\[
\theta^{\prime} = g^{-1}\theta g + g^{-1}dg .
\]
\item \emph{Existence.} Every vector bundle over a manifold admits a connection.
\item \emph{The space of connections.} If $\nabla$ is a connection and $A \in \Omega^1(M, \End E)$, then $\nabla + A$ is a connection, and every connection is $\nabla + A$ for a
unique $A$: the connections form an affine space over $\Omega^1(M, \End E)$.
\end{enumerate}
\end{proposition}

\begin{proof}
(1) Let $s = 0$ on $U$, $p \in U$, and $\chi$ a bump function with $\chi = 1$ near $p$ and support in $U$ (Lemma~\ref{smooth-manifolds:lemma-bump-manifold}). Then $\chi s = 0$, and
$0 = \nabla(\chi s) = d\chi \otimes s + \chi\nabla s$; at $p$, $s(p) = 0$ and $\chi(p) = 1$, so $(\nabla s)_p = 0$. For a section $s$ over $U$ and $p \in U$ define
$(\nabla^Us)_p = (\nabla(\chi s))_p$, where $\chi s$ is extended by $0$ outside $U$ and $\chi$ is as above; by locality this does not depend on $\chi$, it is smooth (near $p$ one
$\chi$ works for all nearby points), and it satisfies the Leibniz rule. It is the only connection on $E|_U$ compatible with restriction, again by locality.

(2) By (1) we may work on $U$. Each $\nabla e_j$ is an $E$-valued $1$-form, so by Proposition~\ref{connections:proposition-bundle-valued-forms}(4) it is uniquely
$\sum_i\theta^i_j \otimes e_i$ with $1$-forms $\theta^i_j$. By linearity and the Leibniz rule,
$\nabla(\sum_js^je_j) = \sum_j(ds^j \otimes e_j + s^j\sum_i\theta^i_j \otimes e_i)$, which is the formula.

(3) On the one hand, by the Leibniz rule,
$\nabla e^{\prime}_j = \nabla(\sum_ke_kg^k_j) = \sum_k(dg^k_j \otimes e_k + g^k_j\sum_i\theta^i_k \otimes e_i) = \sum_i(dg + \theta g)^i_j \otimes e_i$. On the other hand, by definition of
$\theta^{\prime}$, $\nabla e^{\prime}_j = \sum_l\theta^{\prime l}_j \otimes e^{\prime}_l = \sum_i(g\theta^{\prime})^i_j \otimes e_i$. Comparing, $g\theta^{\prime} = dg + \theta g$.

(4) Choose an open cover $(U_\alpha)$ of $M$ by sets over which $E$ has a frame $e_\alpha$, and a partition of unity $(\psi_\alpha)$ subordinate to it
(Theorem~\ref{smooth-manifolds:theorem-partitions-of-unity}). On $U_\alpha$ let $\nabla^\alpha$ be the connection with connection form $0$ in the frame $e_\alpha$, that is,
$\nabla^\alpha(\sum_js^je_{\alpha,j}) = \sum_jds^j \otimes e_{\alpha,j}$. Put $\nabla s = \sum_\alpha\psi_\alpha\nabla^\alpha(s|_{U_\alpha})$, each term extended by $0$ outside the support
of $\psi_\alpha$; the sum is locally finite. It is linear, and
$\nabla(fs) = \sum_\alpha\psi_\alpha(df \otimes s + f\nabla^\alpha s) = df \otimes s + f\nabla s$, because $\sum_\alpha\psi_\alpha = 1$. (A general affine combination of connections
fails the Leibniz rule; the coefficients must add up to $1$.)

(5) $(\nabla + A)(fs) = df \otimes s + f\nabla s + fAs$, so $\nabla + A$ is a connection. If $\nabla^{\prime}$ is another connection, $D = \nabla^{\prime} - \nabla$ satisfies
$D(fs) = fD(s)$: the terms $df \otimes s$ cancel. So $(X, s) \mapsto D(s)(X)$ is $C^\infty(M)$-linear in both arguments, and by Lemma~\ref{connections:lemma-tensoriality} it is given
by a section of $\Hom(TM \otimes E, E) = T^*M \otimes \End E$, that is, by some $A \in \Omega^1(M, \End E)$.
\end{proof}

By (2) and (5), a connection on a trivial bundle $M \times \mathbb{K}^r$ is $d + \theta$ for a matrix $\theta$ of $1$-forms on $M$, its connection form in the standard frame.
The connection form is not intrinsic: the transformation rule (3) shows that the same connection can have connection form $0$ in one frame and $g^{-1}dg \neq 0$ in another. For
instance, on the trivial line bundle over $\RR$ with the connection $d$, the frame $e^{\prime} = e^xe$ has connection form $e^{-x}d(e^x) = dx$. Changing frames is a \emph{gauge
transformation} in the language of physics.

\begin{example}\label{connections:example-connection-forms}
\begin{enumerate}
\item \emph{The sphere.} On $S^2$ minus the poles use the spherical coordinates $x(\vartheta, \varphi) = (\sin\vartheta\cos\varphi, \sin\vartheta\sin\varphi, \cos\vartheta)$ and the
orthonormal frame $e_1 = x_\vartheta = (\cos\vartheta\cos\varphi, \cos\vartheta\sin\varphi, -\sin\vartheta)$, $e_2 = x_\varphi/\sin\vartheta = (-\sin\varphi, \cos\varphi, 0)$ of $TS^2$.
For the projection connection of Example~\ref{connections:example-connections}(4) we differentiate in $\RR^3$ and project onto the tangent plane, which amounts to dropping
the component along the normal $x$. Since $\partial_\vartheta e_1 = -x$, $\partial_\varphi e_1 = \cos\vartheta\,e_2$ and $\partial_\varphi e_2 = -(\cos\varphi, \sin\varphi, 0)
= -(\sin\vartheta\,x + \cos\vartheta\,e_1)$,
\[
de_1 = -x\,d\vartheta + \cos\vartheta\,e_2\,d\varphi, \qquad de_2 = -(\sin\vartheta\,x + \cos\vartheta\,e_1)\,d\varphi,
\]
and dropping the terms along $x$,
\[
\nabla e_1 = \cos\vartheta\,d\varphi \otimes e_2, \qquad \nabla e_2 = -\cos\vartheta\,d\varphi \otimes e_1 .
\]
So the connection form is $\theta = \bigl(\begin{smallmatrix} 0 & -\cos\vartheta\,d\varphi\\ \cos\vartheta\,d\varphi & 0\end{smallmatrix}\bigr)$, a skew-symmetric matrix, in
accordance with Proposition~\ref{connections:proposition-metric-connections} below.
\item \emph{The tautological line bundle of $\CC\PP^1$.} Let $L \subseteq \CC\PP^1 \times \CC^2$ have fibre $L_\ell = \ell$ over the line $\ell$, with the Hermitian inner product of
$\CC^2$ and the projection connection. Over $U_0 = \{[1 : z]\}$ the section $s(z) = (1, z)$ is a frame, $ds = (0, 1)\,dz$, and projecting onto the line spanned by $s$,
\[
\nabla s = \frac{\langle(0, 1), s\rangle}{|s|^2}\,dz \otimes s = \frac{\bar z\,dz}{1 + |z|^2} \otimes s, \qquad \theta = \frac{\bar z\,dz}{1 + |z|^2} .
\]
\end{enumerate}
\end{example}

A connection on $E$ induces connections on all the bundles built from $E$ by linear algebra, by requiring the Leibniz rule for the natural pairings and products.

\begin{proposition}\label{connections:proposition-induced-connections}
Let $\nabla$ and $\nabla^{\prime}$ be connections on $E$ and $E^{\prime}$. The following formulas define connections, uniquely characterised by the Leibniz rules they express.
\begin{enumerate}
\item On $E^\vee$: $(\nabla_X\phi)(s) = X(\phi(s)) - \phi(\nabla_Xs)$. In the dual frame $(e^1, \ldots, e^r)$ of a frame $e$ the connection form is $-\theta^T$.
\item On $E \oplus E^{\prime}$: $\nabla_X(s, s^{\prime}) = (\nabla_Xs, \nabla^{\prime}_Xs^{\prime})$, with connection form $\bigl(\begin{smallmatrix}\theta & 0\\ 0 & \theta^{\prime}\end{smallmatrix}\bigr)$.
\item On $E \otimes E^{\prime}$: $\nabla_X(s \otimes s^{\prime}) = \nabla_Xs \otimes s^{\prime} + s \otimes \nabla^{\prime}_Xs^{\prime}$, with connection form $\theta \otimes 1 + 1 \otimes \theta^{\prime}$
in the frame $(e_i \otimes e^{\prime}_k)$.
\item On $\Hom(E, E^{\prime})$: $(\nabla_XA)(s) = \nabla^{\prime}_X(As) - A(\nabla_Xs)$; in frames, the matrix of $\nabla A$ is $dA + \theta^{\prime}A - A\theta$. On $\End E$ this is
$\nabla A = dA + [\theta, A]$ in matrix form, that is, $\nabla_XA = [\nabla_X, A]$.
\item On $\Lambda^kE$: $\nabla_X(s_1 \wedge \cdots \wedge s_k) = \sum_is_1 \wedge \cdots \wedge \nabla_Xs_i \wedge \cdots \wedge s_k$; on the determinant line bundle $\det E = \Lambda^rE$ the
connection form in the frame $e_1 \wedge \cdots \wedge e_r$ is $\operatorname{tr}\theta$.
\end{enumerate}
\end{proposition}

\begin{proof}
(1) The right side is $C^\infty(M)$-linear in $X$, and in $s$: replacing $s$ by $fs$ gives $X(f\phi(s)) - \phi((Xf)s + f\nabla_Xs) = f(X(\phi(s)) - \phi(\nabla_Xs))$. By
Lemma~\ref{connections:lemma-tensoriality} it defines a section $\nabla_X\phi$ of $E^\vee$, and the three properties of
Proposition~\ref{connections:proposition-connection-operators} are checked directly; for instance $(\nabla_X(f\phi))(s) = X(f\phi(s)) - f\phi(\nabla_Xs) = (Xf)\phi(s) + f(\nabla_X\phi)(s)$.
In the dual frame, $(\nabla e^j)(e_k) = d(\delta^j_k) - e^j(\nabla e_k) = -\theta^j_k$, so $\nabla e^j = -\sum_k\theta^j_k \otimes e^k$, whose matrix is $-\theta^T$. (2) is clear.

(3) Define $\nabla$ on $E \otimes E^{\prime}$ over a set with frames $e$, $e^{\prime}$ by the connection form $\theta \otimes 1 + 1 \otimes \theta^{\prime}$, that is, by
$\nabla(e_i \otimes e^{\prime}_k) = \nabla e_i \otimes e^{\prime}_k + e_i \otimes \nabla^{\prime}e^{\prime}_k$ and the Leibniz rule. For sections $s = \sum_is^ie_i$ and
$s^{\prime} = \sum_ks^{\prime k}e^{\prime}_k$, expanding $\nabla(\sum_{i,k}s^is^{\prime k}e_i \otimes e^{\prime}_k)$ with $d(s^is^{\prime k}) = ds^i\,s^{\prime k} + s^i\,ds^{\prime k}$ gives
$\nabla s \otimes s^{\prime} + s \otimes \nabla^{\prime}s^{\prime}$. So the product rule holds, and since it determines $\nabla$ on the products of frame sections, which span the fibres,
the local definitions do not depend on the frames and glue.

(4) $\Hom(E, E^{\prime}) = E^\vee \otimes E^{\prime}$, and the connection of (1) and (3) gives the stated formula: for $A = \phi \otimes s^{\prime}$,
$\nabla_X(A)(s) = (\nabla_X\phi)(s)s^{\prime} + \phi(s)\nabla^{\prime}_Xs^{\prime} = \nabla^{\prime}_X(\phi(s)s^{\prime}) - \phi(\nabla_Xs)s^{\prime}$. In frames, $A$ is a matrix of functions,
$As = A\sigma$ for the component vector $\sigma$ of $s$, and $\nabla^{\prime}(A\sigma) - A(d\sigma + \theta\sigma) = (dA)\sigma + A\,d\sigma + \theta^{\prime}A\sigma - A\,d\sigma - A\theta\sigma$.

(5) As in (3), define $\nabla$ on $\Lambda^kE$ in the frame $(e_{i_1} \wedge \cdots \wedge e_{i_k})$ by the stated formula on the frame elements and the Leibniz rule; expanding
$s_1 \wedge \cdots \wedge s_k$ in the frame shows that the formula holds for all sections, so the definition does not depend on the frame. For $k = r$,
$\nabla(e_1 \wedge \cdots \wedge e_r) = \sum_ie_1 \wedge \cdots \wedge (\sum_l\theta^l_ie_l) \wedge \cdots \wedge e_r = (\sum_i\theta^i_i)\,e_1 \wedge \cdots \wedge e_r$, since the terms with
$l \neq i$ contain a repeated factor.
\end{proof}

\section{Curvature}\label{connections:section-curvature}

For the trivial connection on a trivial bundle, second derivatives commute: $X(Y(s)) - Y(X(s)) = [X, Y](s)$, the Lie bracket measuring exactly the failure of $X$ and $Y$ to commute
as derivations. For a general connection the difference $\nabla_X\nabla_Ys - \nabla_Y\nabla_Xs - \nabla_{[X, Y]}s$ need not vanish, and it turns out to be a tensor: the curvature.
The cleanest way to see this is to extend $\nabla$ to bundle-valued forms, as $d$ extends from functions to forms, and to observe that the extension does not square to zero.

\begin{proposition}[Exterior covariant derivative]\label{connections:proposition-exterior-covariant}
Let $\nabla$ be a connection on $E$. There is a unique family of $\mathbb{K}$-linear maps $d_\nabla \colon \Omega^k(M, E) \to \Omega^{k+1}(M, E)$, $k \geq 0$, such that
$d_\nabla = \nabla$ on $\Omega^0(M, E) = \Gamma(E)$ and
\[
d_\nabla(\alpha \wedge \omega) = d\alpha \wedge \omega + (-1)^k\alpha \wedge d_\nabla\omega \qquad (\alpha \in \Omega^k(M),\ \omega \in \Omega^l(M, E)).
\]
In a frame with connection form $\theta$, writing $\omega = \sum_i\omega^i \otimes e_i$,
\[
d_\nabla\omega = \sum_i\Bigl(d\omega^i + \sum_j\theta^i_j \wedge \omega^j\Bigr) \otimes e_i, \qquad\text{in vector notation}\quad d_\nabla\omega = d\omega + \theta \wedge \omega .
\]
On vector fields, for $\omega \in \Omega^k(M, E)$,
\[
(d_\nabla\omega)(X_0, \ldots, X_k) = \sum_i(-1)^i\nabla_{X_i}\bigl(\omega(X_0, \ldots, \widehat{X_i}, \ldots, X_k)\bigr) + \sum_{i<j}(-1)^{i+j}\omega([X_i, X_j], X_0, \ldots, \widehat{X_i}, \ldots, \widehat{X_j}, \ldots, X_k);
\]
for $k = 1$, $(d_\nabla\omega)(X, Y) = \nabla_X(\omega(Y)) - \nabla_Y(\omega(X)) - \omega([X, Y])$.
\end{proposition}

\begin{proof}
\emph{Uniqueness.} As for the exterior derivative (Theorem~\ref{differential-forms:theorem-exterior-derivative}), the Leibniz rule makes $d_\nabla$ local: if $\omega = 0$ near $p$,
choose a bump $\chi$ equal to $1$ near $p$ with $\chi\omega = 0$; then $0 = d_\nabla(\chi\omega) = d\chi \wedge \omega + \chi d_\nabla\omega$, and at $p$ this is $(d_\nabla\omega)_p$. So it
suffices to compute in a frame, where $\omega = \sum_i\omega^i \wedge e_i$ with the $e_i$ viewed as $E$-valued $0$-forms. The Leibniz rule forces
\[
d_\nabla\omega = \sum_i\bigl(d\omega^i \wedge e_i + (-1)^k\omega^i \wedge \nabla e_i\bigr) = \sum_id\omega^i\,e_i + \sum_{i,j}(-1)^k\omega^i \wedge \theta^j_i\,e_j
= \sum_j\Bigl(d\omega^j + \sum_i\theta^j_i \wedge \omega^i\Bigr)e_j,
\]
using $\omega^i \wedge \theta^j_i = (-1)^k\theta^j_i \wedge \omega^i$ for the $k$-form $\omega^i$ and the $1$-form $\theta^j_i$.

\emph{Existence.} Define $d_\nabla$ on a frame domain by the formula $d\omega + \theta \wedge \omega$. It does not depend on the frame: if $e^{\prime} = eg$, then $\omega = g\omega^{\prime}$ and
$\theta^{\prime} = g^{-1}\theta g + g^{-1}dg$ (Proposition~\ref{connections:proposition-connections-exist}(3)), and
\[
g(d\omega^{\prime} + \theta^{\prime} \wedge \omega^{\prime}) = g\,d\omega^{\prime} + \theta \wedge g\omega^{\prime} + dg \wedge \omega^{\prime} = d(g\omega^{\prime}) + \theta \wedge (g\omega^{\prime}) = d\omega + \theta \wedge \omega,
\]
which says that the two definitions give the same $E$-valued form. So the local definitions glue. For $k = 0$ the formula is that of $\nabla$
(Proposition~\ref{connections:proposition-connections-exist}(2)), and the Leibniz rule holds in a frame: $d(\alpha\omega^i) + \sum_j\theta^i_j \wedge \alpha\omega^j
= d\alpha \wedge \omega^i + (-1)^k\alpha \wedge (d\omega^i + \sum_j\theta^i_j \wedge \omega^j)$, since $\theta^i_j \wedge \alpha = (-1)^k\alpha \wedge \theta^i_j$.

\emph{The formula with vector fields.} Both sides are local and $\mathbb{K}$-linear in $\omega$, and locally $\omega$ is a sum of forms $\beta \otimes s$ with $\beta \in \Omega^k$ and $s$ a
section. For such a form, the left side is $(d\beta \otimes s + (-1)^k\beta \wedge \nabla s)(X_0, \ldots, X_k)$. On the right side,
$\nabla_{X_i}(\beta(\ldots)s) = X_i(\beta(\ldots))s + \beta(\ldots)\nabla_{X_i}s$; the terms of the first kind, together with the bracket terms, give $d\beta(X_0, \ldots, X_k)\,s$ by the
formula of Theorem~\ref{differential-forms:theorem-exterior-derivative}. The remaining terms $\sum_i(-1)^i\beta(X_0, \ldots, \widehat{X_i}, \ldots, X_k)\nabla_{X_i}s$ equal
$(-1)^k(\beta \wedge \nabla s)(X_0, \ldots, X_k)$: by Proposition~\ref{multilinear-algebra:proposition-wedge-formula}, the shuffle that moves $X_i$ to the last place has sign
$(-1)^{k-i}$, and $(-1)^k(-1)^{k-i} = (-1)^i$.
\end{proof}

\begin{definition}\label{connections:definition-curvature}
The \emph{exterior covariant derivative} of $\nabla$ is the operator $d_\nabla$ of Proposition~\ref{connections:proposition-exterior-covariant}; it is determined by
$d_\nabla(\alpha \otimes s) = d\alpha \otimes s + (-1)^k\alpha \wedge \nabla s$ for $\alpha \in \Omega^k(M)$ and $s \in \Gamma(E)$. The \emph{curvature} of $\nabla$ is
$F_\nabla = d_\nabla \circ \nabla \colon \Gamma(M, E) \to \Omega^2(M, E)$, also written $F$ or $R$.
\end{definition}

\begin{proposition}\label{connections:proposition-curvature}
$F_\nabla$ is $C^\infty(M)$-linear, hence an element of $\Omega^2(M, \End E)$, and
\[
F_\nabla(X, Y)s = \nabla_X\nabla_Ys - \nabla_Y\nabla_Xs - \nabla_{[X, Y]}s .
\]
In a frame with connection form $\theta$ one has $F = d\theta + \theta \wedge \theta$, and under a change of frame $F' = g^{-1}Fg$. Moreover $d_\nabla F = 0$ in the sense that
$dF + \theta \wedge F - F \wedge \theta = 0$ (the \emph{Bianchi identity}), and $d_\nabla \circ d_\nabla = F_\nabla \wedge -$ on $\Omega^\bullet(M, E)$.
\end{proposition}

\begin{proof}
\emph{Tensoriality.} Using the Leibniz rules of $\nabla$ and $d_\nabla$ and $d(df) = 0$,
\[
F(fs) = d_\nabla(df \otimes s + f\nabla s) = \bigl(d(df) \otimes s - df \wedge \nabla s\bigr) + \bigl(df \wedge \nabla s + f\,d_\nabla\nabla s\bigr) = fF(s).
\]
So $(X, Y, s) \mapsto F(s)(X, Y)$ is $C^\infty(M)$-linear in all three arguments, and by Lemma~\ref{connections:lemma-tensoriality} it is given by an element of
$\Omega^2(M, \End E)$: at each point, $F(v, w)$ is an endomorphism of $E_p$, alternating in $v$ and $w$.

\emph{The formula with vector fields} is the case $k = 1$ of Proposition~\ref{connections:proposition-exterior-covariant} applied to $\omega = \nabla s$:
$F(s)(X, Y) = \nabla_X((\nabla s)(Y)) - \nabla_Y((\nabla s)(X)) - (\nabla s)([X, Y])$.

\emph{In a frame,} $Fe_j = d_\nabla(\sum_i\theta^i_j \otimes e_i) = \sum_i(d\theta^i_j + \sum_k\theta^i_k \wedge \theta^k_j) \otimes e_i$ by the frame formula for $d_\nabla$, so the matrix of $F$
(with $Fe_j = \sum_iF^i_je_i$) is $d\theta + \theta \wedge \theta$. Since $F$ is a tensor, its matrices in two frames are related as the matrices of any endomorphism:
$Fe^{\prime}_j = \sum_kg^k_jFe_k$, so $F^{\prime} = g^{-1}Fg$. (This can also be checked from the transformation rule of $\theta$.)

\emph{Bianchi.} For matrices of $1$-forms, $d(\theta \wedge \theta) = d\theta \wedge \theta - \theta \wedge d\theta$, so
$dF = d\theta \wedge \theta - \theta \wedge d\theta = (F - \theta \wedge \theta) \wedge \theta - \theta \wedge (F - \theta \wedge \theta) = F \wedge \theta - \theta \wedge F$. The induced
connection on $\End E$ (Proposition~\ref{connections:proposition-induced-connections}(4)) has exterior derivative $d_\nabla A = dA + \theta \wedge A - (-1)^kA \wedge \theta$ on
$\End E$-valued $k$-forms, by the same argument as in Proposition~\ref{connections:proposition-exterior-covariant}; for $k = 2$ the Bianchi identity says $d_\nabla F = 0$.

\emph{$d_\nabla^2$.} In a frame, $d_\nabla d_\nabla\omega = d(d\omega + \theta \wedge \omega) + \theta \wedge (d\omega + \theta \wedge \omega) = d\theta \wedge \omega - \theta \wedge d\omega + \theta \wedge d\omega
+ \theta \wedge \theta \wedge \omega = F \wedge \omega$.
\end{proof}

So $d_\nabla$ is a ``derivative'' on bundle-valued forms that squares to zero exactly when the curvature vanishes; the de Rham complex of a flat bundle (Chapter~\ref{de-rham})
is built from it. The term $\nabla_{[X, Y]}$ in the curvature is what makes it a tensor: for coordinate vector fields, $[\partial_i, \partial_j] = 0$ and
$F(\partial_i, \partial_j) = \nabla_i\nabla_j - \nabla_j\nabla_i$ is the commutator of covariant partial derivatives; in coordinates
$F(\partial_i, \partial_j) = \partial_i\theta_j - \partial_j\theta_i + [\theta_i, \theta_j]$, where $\theta = \sum_i\theta_i\,dx^i$ with matrix-valued functions $\theta_i$.

\begin{example}\label{connections:example-curvature-computations}
\begin{enumerate}
\item \emph{Trivial bundles and line bundles.} The trivial connection $d$ has $F = 0$. The connection $d + A$ on $M \times \mathbb{K}^r$ has $F = dA + A \wedge A$; for a line bundle
the connection form is a single $1$-form $A$, and $F = dA$. On a manifold of dimension $1$ every connection is flat, since there are no nonzero $2$-forms.
\item \emph{The sphere.} With the connection form of Example~\ref{connections:example-connection-forms}(1), all entries of $\theta$ are multiples of $d\varphi$, so $\theta \wedge \theta = 0$, and
$d(\cos\vartheta\,d\varphi) = -\sin\vartheta\,d\vartheta \wedge d\varphi$. With the area form $dA = \sin\vartheta\,d\vartheta \wedge d\varphi$ (for which $dA(e_1, e_2) = 1$),
\[
F = d\theta = \begin{pmatrix} 0 & dA\\ -dA & 0\end{pmatrix}, \qquad F(e_1, e_2)e_1 = -e_2, \quad F(e_1, e_2)e_2 = e_1 .
\]
So $F = dA \otimes J$, where $J$ is the rotation of each tangent plane by $-\pi/2$ ($Je_1 = -e_2$, $Je_2 = e_1$). Although $\theta$ was only defined away from the poles, $F$ is
a global tensor, and $dA \otimes J$ is indeed defined on all of $S^2$. The number $\langle F(e_1, e_2)e_2, e_1\rangle = 1$ is the Gaussian curvature of the unit sphere
(Chapter~\ref{riemannian-geometry}); since $F \neq 0$, no neighbourhood of any point of $S^2$ has a frame of parallel vector fields.
\item \emph{The tautological line bundle of $\CC\PP^1$.} With $\theta = \bar z\,dz/(1 + |z|^2)$ from Example~\ref{connections:example-connection-forms}(2), and writing $z = x + iy$,
\[
F = d\theta = \frac{\partial}{\partial\bar z}\Bigl(\frac{\bar z}{1 + |z|^2}\Bigr)\,d\bar z \wedge dz = \frac{d\bar z \wedge dz}{(1 + |z|^2)^2} = \frac{2i\,dx \wedge dy}{(1 + x^2 + y^2)^2},
\]
using $\partial_{\bar z}(\bar z(1 + z\bar z)^{-1}) = (1 + |z|^2)^{-2}$ and $d\bar z \wedge dz = 2i\,dx \wedge dy$. The real $2$-form $\frac{i}{2\pi}F = -\frac1\pi\,dx \wedge dy/(1 + x^2 + y^2)^2$
integrates over $U_0 \cong \CC$, which is $\CC\PP^1$ minus a point, to $-\frac1\pi\int_0^\infty\frac{2\pi r\,dr}{(1 + r^2)^2} = -1$. This integer is the first Chern class of the
tautological bundle (Chapter~\ref{characteristic-classes}); that the integral of a curvature form is an integer independent of the connection is the beginning of Chern--Weil
theory.
\end{enumerate}
\end{example}

\begin{proposition}\label{connections:proposition-pullback-tensor}
Let $\nabla$ be a connection on $E \to M$ with curvature $F$, and $\nabla'$ a connection on $E' \to M$ with curvature $F'$.
\begin{enumerate}
\item For smooth $f \colon N \to M$ there is a unique connection $f^*\nabla$ on $f^*E$ with $(f^*\nabla)(f^*s) = f^*(\nabla s)$ for sections $s$ of $E$; its curvature is $f^*F$.
\item The connection $\nabla \otimes 1 + 1 \otimes \nabla'$ on $E \otimes E'$ has curvature $F \otimes 1 + 1 \otimes F'$, and the dual connection on $E^\vee$, $(\nabla^\vee\phi)(s) = d(\phi(s)) - \phi(\nabla s)$, has curvature $-F^T$.
\item For a line bundle, $\nabla$ and $\nabla + \alpha$, with $\alpha$ a $1$-form, have curvatures differing by $d\alpha$; so the curvature of a line bundle is a closed $2$-form whose class does not depend on the connection.
\end{enumerate}
\end{proposition}

Here $f^*s = s \circ f$ is the pulled back section (Proposition~\ref{vector-bundles:proposition-pullback-fibre-product}), and $f^*(\nabla s)$ is the $f^*E$-valued $1$-form
$v \mapsto (\nabla s)(df(v))$ on $N$. Concretely, $(f^*\nabla)_v(s \circ f) = \nabla_{df(v)}s$ for $v \in T_qN$: the pullback connection differentiates in the direction of the image
vector. In a frame $e$ of $E$ with connection form $\theta$, the pulled back frame $e \circ f$ has connection form $f^*\theta$.

\begin{proof}
(1) Over a frame domain $U$ of $E$ with connection form $\theta$, define a connection on $f^*E|_{f^{-1}U}$ by the connection form $f^*\theta$ in the frame $(e_j \circ f)$. If
$e^{\prime} = eg$ is another frame, the pulled back frames are related by $g \circ f$, and pulling back the transformation rule gives
$f^*\theta^{\prime} = (g \circ f)^{-1}(f^*\theta)(g \circ f) + (g \circ f)^{-1}d(g \circ f)$, since $f^*$ commutes with $d$ and with products. So the local connections agree on overlaps
and glue. For $s = \sum_js^je_j$, $f^*s = \sum_j(s^j \circ f)(e_j \circ f)$, and $(f^*\nabla)(f^*s) = \sum_i(d(s^i \circ f) + \sum_j(f^*\theta^i_j)(s^j \circ f)) \otimes (e_i \circ f)
= f^*(\nabla s)$. Uniqueness holds because the condition fixes $f^*\nabla$ on the pulled back frames, which span $f^*E$ locally over $C^\infty(N)$. The curvature, computed in the
pulled back frame, is $d(f^*\theta) + f^*\theta \wedge f^*\theta = f^*(d\theta + \theta \wedge \theta) = f^*F$.

(2) In frames $(e_j)$, $(e'_k)$ with connection forms $\theta$, $\theta'$, the frame $(e_j \otimes e'_k)$ has connection form $\theta \otimes 1 + 1 \otimes \theta'$
(Proposition~\ref{connections:proposition-induced-connections}), and as the two summands act on different factors, they commute:
$(\theta \otimes 1) \wedge (1 \otimes \theta^{\prime}) + (1 \otimes \theta^{\prime}) \wedge (\theta \otimes 1) = 0$ by graded commutativity of the $1$-form entries. So $d$ of the connection form plus its
square is $(d\theta + \theta \wedge \theta) \otimes 1 + 1 \otimes (d\theta' + \theta' \wedge \theta')$. The dual frame has connection form $-\theta^T$, and
$d(-\theta^T) + \theta^T \wedge \theta^T = -(d\theta + \theta \wedge \theta)^T$, as $(\theta \wedge \theta)^T = -\theta^T \wedge \theta^T$ for matrices of $1$-forms.

(3) For rank one, $\theta$ is a $1$-form and $\theta \wedge \theta = 0$, so $F = d\theta$ in every frame, and replacing $\theta$ by $\theta + \alpha$ adds $d\alpha$. Closedness is $dF = F \wedge \theta - \theta \wedge F = 0$, the
Bianchi identity for commuting $1 \times 1$ matrices.
\end{proof}

\section{Parallel transport, holonomy and flatness}\label{connections:section-parallel-transport}

A connection lets us say when a section is ``constant'' along a curve: when its covariant derivative in the direction of the curve vanishes. Solving this differential
equation moves vectors from one fibre to another along the curve. First we make precise what the derivative along a curve means.

\begin{remark}\label{connections:remark-along-curves}
A connection differentiates sections of $E$ over $M$ in the direction of vector fields on $M$, while along a curve $\gamma \colon I \to M$ the direction $\dot\gamma$ and
the section to be differentiated are only defined along the curve. The notation $\nabla_{\dot\gamma}s$ is justified as follows.
\begin{enumerate}
\item \emph{The direction may be a single vector.} For a section $s$ of $E$, $\nabla s$ is an $E$-valued $1$-form, so $\nabla_vs = (\nabla s)_p(v) \in E_p$ is defined for
every tangent vector $v \in T_pM$; equivalently, $\nabla_Xs$ at $p$ depends only on $X_p$, because $\nabla_{fX}s = f\nabla_Xs$. So for a section $s$ defined near
$\gamma(I)$, $\nabla_{\dot\gamma(t)}s$ makes sense without extending $\dot\gamma$ to a vector field.
\item \emph{Only the values of $s$ along a curve matter.} In a local frame, $\nabla_vs = \sum_i\bigl(v(s^i) + \sum_j\theta^i_j(v)s^j\bigr)e_i$
(Proposition~\ref{connections:proposition-connections-exist}), and $v(s^i) = (s^i \circ c)^{\prime}(0)$ for any curve $c$ with $\dot c(0) = v$. So $\nabla_vs$ depends
only on the values of $s$ along such a curve.
\item \emph{Sections along a curve need not extend.} A \emph{section along $\gamma$} is a section of $\gamma^*E$, that is, a smooth map $t \mapsto s(t) \in E_{\gamma(t)}$;
for instance the velocity $\dot\gamma$ is a section of $\gamma^*TM$. It need not be of the form $S \circ \gamma$ for a section $S$ of $E$, even near one point. For
$\gamma(t) = (t^2, 0)$ in $\RR^2$ we have $\gamma(t) = \gamma(-t)$ but $\dot\gamma(t) = (2t, 0) \neq \dot\gamma(-t)$ for $t \neq 0$, so no vector field $X$ near the origin
satisfies $X(\gamma(t)) = \dot\gamma(t)$; the same happens at every self-intersection of a curve, where two velocities sit at one point.
\end{enumerate}
The precise definition therefore uses the pullback connection of Proposition~\ref{connections:proposition-pullback-tensor}(1): for a section $s$ along $\gamma$,
\[
\nabla_{\dot\gamma}s = (\gamma^*\nabla)_{d/dt}\,s ,
\]
a derivative on the interval $I$, where $d/dt$ is a genuine vector field. By the defining property $(\gamma^*\nabla)(\gamma^*S) = \gamma^*(\nabla S)$,
$(\gamma^*\nabla)_{d/dt}(S \circ \gamma) = (\nabla S)(\dot\gamma) = \nabla_{\dot\gamma}S$ in the sense of (1), and in a local frame
$(\gamma^*\nabla)_{d/dt}s = \sum_i\bigl(\dot s^i + \sum_j\theta^i_j(\dot\gamma)s^j\bigr)e_i$. When $s$ does extend to a section $S$ near $\gamma(t)$, the two meanings of
$\nabla_{\dot\gamma}s$ agree, whatever the extension, by (2).
\end{remark}

\begin{definition}\label{connections:definition-parallel}
A section $s$ along a smooth curve $\gamma \colon [0, 1] \to M$, that is, a section of $\gamma^*E$, is \emph{parallel} if $\nabla_{\dot\gamma}s = 0$, in the sense of
Remark~\ref{connections:remark-along-curves}. \emph{Parallel transport} along $\gamma$ is the map $P_\gamma \colon E_{\gamma(0)} \to E_{\gamma(1)}$
sending $v$ to $s(1)$, where $s$ is the parallel section with $s(0) = v$. For piecewise smooth curves, transport along the smooth pieces is composed.
\end{definition}

\begin{proposition}\label{connections:proposition-parallel-transport}
Parallel sections along $\gamma$ exist and are unique for each initial value, so $P_\gamma$ is a well-defined linear isomorphism, depending smoothly on the data and satisfying
$P_{\gamma \ast \delta} = P_\delta \circ P_\gamma$ and $P_{\bar\gamma} = P_\gamma^{-1}$.
\end{proposition}

\begin{proof}
In a local frame the equation $\nabla_{\dot\gamma}s = 0$ reads
\[
\dot s^i(t) + \sum_j\theta^i_j(\dot\gamma(t))\,s^j(t) = 0, \qquad i = 1, \ldots, r,
\]
a linear system of ordinary differential equations with smooth coefficients. By Proposition~\ref{real-analysis:proposition-ode-dependence}(2) its solutions exist on the whole
interval on which $\gamma$ stays in the frame domain, are unique, and depend linearly on the initial value, and smoothly on parameters
(Theorem~\ref{real-analysis:theorem-smooth-dependence}). The interval $[0, 1]$ is covered by finitely many subintervals on each of which $\gamma$ stays in one frame domain; solving
successively on them, using the end value of one piece as initial value of the next, gives a unique parallel section on $[0, 1]$. Uniqueness does not depend on the frames,
because the equation $\nabla_{\dot\gamma}s = 0$ does not. So $P_\gamma$ is well defined and linear. Running the curve backwards, $\bar\gamma(t) = \gamma(1 - t)$, the section
$t \mapsto s(1 - t)$ is parallel along $\bar\gamma$ (the equation is linear in $\dot\gamma$), so $P_{\bar\gamma}P_\gamma = \id$ and $P_\gamma$ is invertible; the composition rule follows
from uniqueness. The same argument shows that $P_\gamma$ does not change under an increasing reparametrisation of $\gamma$.
\end{proof}

\begin{proposition}[Homotopy invariance of pullbacks]\label{connections:proposition-homotopy-invariance}
Let $E \to N$ be a vector bundle and $F \colon M \times \RR \to N$ a smooth map, and put $f = F(-, 0)$ and $g = F(-, 1)$. Then $f^*E \cong g^*E$ as vector bundles on $M$. In particular every vector bundle on a star-shaped open
subset of $\RR^n$, for instance on $\RR^n$ or on a ball, is trivial.
\end{proposition}

\begin{proof}
Choose a connection on $E' = F^*E$ (Proposition~\ref{connections:proposition-connections-exist}). For $x \in M$ let $\gamma_x(t) = (x, t)$, $0 \leq t \leq 1$, and let $P_x \colon E'_{(x,0)} \to E'_{(x,1)}$ be parallel
transport along $\gamma_x$. Since $E'_{(x,0)} = E_{f(x)} = (f^*E)_x$ and $E'_{(x,1)} = (g^*E)_x$, the $P_x$ are linear isomorphisms $(f^*E)_x \to (g^*E)_x$, and they depend smoothly on $x$
(Proposition~\ref{connections:proposition-parallel-transport}), so they form an isomorphism of vector bundles $f^*E \to g^*E$. For a star-shaped $U$, with centre $p$, apply this to $F(x, t) = p + t(x - p)$ and
$N = U$: then $g = \id$ and $f$ is constant, so $E = \id^*E \cong f^*E = U \times E_p$ is trivial.
\end{proof}

\begin{example}\label{connections:example-transport}
\begin{enumerate}
\item For the trivial connection $d$ on $M \times \mathbb{K}^r$ the parallel sections along any curve are the constant ones, and $P_\gamma = \id$.
\item \emph{Around a latitude of the sphere.} Let $c(\varphi) = x(\vartheta_0, \varphi)$, $\varphi \in [0, 2\pi]$, be the circle of latitude at polar angle $\vartheta_0$, with the frame
and connection form of Example~\ref{connections:example-connection-forms}(1). For $V = ae_1 + be_2$, $\nabla_{\dot c}e_1 = \cos\vartheta_0\,e_2$ and
$\nabla_{\dot c}e_2 = -\cos\vartheta_0\,e_1$, so $\nabla_{\dot c}V = (a^{\prime} - b\cos\vartheta_0)e_1 + (b^{\prime} + a\cos\vartheta_0)e_2$. With $\kappa = \cos\vartheta_0$ the parallel
fields are
\[
a(\varphi) = a_0\cos(\kappa\varphi) + b_0\sin(\kappa\varphi), \qquad b(\varphi) = -a_0\sin(\kappa\varphi) + b_0\cos(\kappa\varphi):
\]
relative to the frame $(e_1, e_2)$, which itself comes back to its initial value, the vector rotates by the angle $-\kappa\varphi$. After one turn it is rotated by
$-2\pi\cos\vartheta_0$, which modulo $2\pi$ equals $2\pi(1 - \cos\vartheta_0)$, the area of the spherical cap $\{\vartheta < \vartheta_0\}$ enclosed by the circle. On the equator,
$\vartheta_0 = \pi/2$, the transport is trivial, and near the north pole the rotation is small, as the enclosed area is. This is the Gauss--Bonnet theorem in its simplest form
(Chapter~\ref{riemannian-geometry}), and an instance of the next proposition.
\end{enumerate}
\end{example}

The next result is the geometric meaning of curvature: transport around a small loop differs from the identity by the curvature times the enclosed area.

\begin{proposition}[Holonomy of small loops]\label{connections:proposition-small-loops}
Let $x^1, \ldots, x^n$ be coordinates centred at $p$, and for small $\varepsilon > 0$ let $\gamma_\varepsilon$ be the loop at $p$ that runs around the square with vertices $0$,
$\varepsilon e_1$, $\varepsilon(e_1 + e_2)$, $\varepsilon e_2$ in the $(x^1, x^2)$-plane, in this order. Then
\[
P_{\gamma_\varepsilon} = \id - \varepsilon^2F(\partial_1, \partial_2)_p + O(\varepsilon^3).
\]
\end{proposition}

\begin{proof}
Work in a frame near $p$, with $\theta = \sum_i\theta_i\,dx^i$, and write $\theta_i$, $\partial_j\theta_i$ for the values at $p$. Transport along a segment $t \mapsto q + tu$, $u$ a unit
coordinate vector $\pm e_k$, solves $\dot s = -\theta(u)s$ with $\theta(u) = \pm\theta_k(q + tu)$. Differentiating the equation, $\ddot s = -(\partial_u\theta(u))s + \theta(u)^2s$, so by
Taylor's formula (the solution depends smoothly on $q$ and $\varepsilon$) transport along the segment of length $\varepsilon$ is
$T = \id - \varepsilon\theta(u)(q) + \frac{\varepsilon^2}2\bigl(\theta(u)^2 - \partial_u\theta(u)\bigr)(q) + O(\varepsilon^3)$. For the four sides, expanding the values at the starting
points $0$, $\varepsilon e_1$, $\varepsilon(e_1 + e_2)$, $\varepsilon e_2$ around $p$ to first order,
\[
\begin{aligned}
T_1 &= \id - \varepsilon\theta_1 + \tfrac{\varepsilon^2}2(\theta_1^2 - \partial_1\theta_1),\\
T_2 &= \id - \varepsilon\theta_2 - \varepsilon^2\partial_1\theta_2 + \tfrac{\varepsilon^2}2(\theta_2^2 - \partial_2\theta_2),\\
T_3 &= \id + \varepsilon\theta_1 + \varepsilon^2(\partial_1\theta_1 + \partial_2\theta_1) + \tfrac{\varepsilon^2}2(\theta_1^2 - \partial_1\theta_1),\\
T_4 &= \id + \varepsilon\theta_2 + \varepsilon^2\partial_2\theta_2 + \tfrac{\varepsilon^2}2(\theta_2^2 - \partial_2\theta_2),
\end{aligned}
\]
up to $O(\varepsilon^3)$; for $T_3$ and $T_4$ the direction is $-e_1$, $-e_2$, so $\theta(u) = -\theta_1$, $-\theta_2$ and $\partial_u\theta(u) = \partial_1\theta_1$, $\partial_2\theta_2$. In
$P_{\gamma_\varepsilon} = T_4T_3T_2T_1$ the terms of order $\varepsilon$ cancel. The terms of order $\varepsilon^2$ coming from a single factor add up to
$\theta_1^2 + \theta_2^2 - \partial_1\theta_2 + \partial_2\theta_1$, and the products of two first-order terms, taken in the order of the factors, add up to
$\theta_2\theta_1 - \theta_2^2 - \theta_2\theta_1 - \theta_1\theta_2 - \theta_1^2 + \theta_2\theta_1 = -\theta_1^2 - \theta_2^2 - \theta_1\theta_2 + \theta_2\theta_1$. The total coefficient
of $\varepsilon^2$ is $-(\partial_1\theta_2 - \partial_2\theta_1 + \theta_1\theta_2 - \theta_2\theta_1) = -F(\partial_1, \partial_2)_p$, by the coordinate formula for $F$ after
Proposition~\ref{connections:proposition-curvature}.
\end{proof}

For a line bundle the formula is exact and global in the following sense.

\begin{proposition}\label{connections:proposition-line-holonomy}
Let $L$ be a line bundle with connection $\nabla$ and curvature $F$ (a closed $2$-form), let $\bar D$ be the closed unit disc, $h \colon B \to M$ a smooth map defined on an
open disc $B \supseteq \bar D$, and let $\gamma$ be the loop $h|_{\partial\bar D}$, traversed counterclockwise. Then parallel transport around $\gamma$ is multiplication by
\[
P_\gamma = \exp\Bigl(-\int_{\bar D}h^*F\Bigr).
\]
\end{proposition}

\begin{proof}
The bundle $h^*L$ over the star-shaped $B$ is trivial (Proposition~\ref{connections:proposition-homotopy-invariance}), and
parallel transport along $\gamma$ for $\nabla$ is parallel transport along the boundary circle for $h^*\nabla$, because $(h \circ c)^*\nabla = c^*(h^*\nabla)$ for curves $c$ in the
disc (both satisfy the characterisation of Proposition~\ref{connections:proposition-pullback-tensor}(1)). In a frame over $B$, $h^*\nabla = d + \alpha$ for a $1$-form $\alpha$, and a
parallel section along the boundary circle $c$ solves the scalar equation $\dot s + \alpha(\dot c)s = 0$, so $s(1) = \exp(-\int_c\alpha)\,s(0)$. By Stokes' theorem
(Theorem~\ref{differential-forms:theorem-stokes}), $\int_c\alpha = \int_{\bar D}d\alpha = \int_{\bar D}h^*F$, as $h^*F = d\alpha$ by
Proposition~\ref{connections:proposition-pullback-tensor}(1),(3).
\end{proof}

\begin{definition}\label{connections:definition-holonomy}
The \emph{holonomy group} of $\nabla$ at $p$ is $\mathrm{Hol}_p(\nabla) = \{P_\gamma : \gamma \text{ a loop at } p\} \subseteq \GL(E_p)$, a subgroup. A connection is \emph{flat} if $F_\nabla = 0$.
\end{definition}

It is a subgroup by Proposition~\ref{connections:proposition-parallel-transport}: the product of $P_\gamma$ and $P_\delta$ is transport along the concatenated loop, and the inverse is
transport along the reversed loop. If $\sigma$ is a path from $p$ to $q$, then $\mathrm{Hol}_q = P_\sigma\mathrm{Hol}_pP_\sigma^{-1}$, so on a connected manifold the holonomy groups at
different points are conjugate. For the sphere, Example~\ref{connections:example-transport}(2) shows that the holonomy contains rotations by all angles $2\pi(1 - \cos\vartheta_0)$,
hence it is the whole rotation group $\mathrm{SO}(2)$ of the tangent plane.

\begin{theorem}\label{connections:theorem-flat-local-systems}
For a connection $\nabla$ on $E \to M$ the following are equivalent:
\begin{enumerate}
\item $F_\nabla = 0$;
\item every point has a neighbourhood with a frame of parallel sections (a \emph{flat frame}), in which $\nabla = d$;
\item parallel transport along a path depends only on its homotopy class with fixed endpoints.
\end{enumerate}
In that case the sheaf $\mathcal{E}^\nabla$ of parallel sections is a local system of rank $r$ with $\mathcal{E}^\nabla \otimes_{\RR_M} C^\infty_M \cong \mathcal{E}$, and on a connected $M$ this gives an
equivalence between flat bundles and representations of $\pi_1(M, p)$, the representation being the holonomy
\[
\pi_1(M, p) \to \GL(E_p), \qquad [\gamma] \mapsto P_\gamma^{-1}.
\]
\end{theorem}

\begin{proof}
(2)$\Rightarrow$(1) is clear, since $d$ has zero curvature, and curvature is local.

(1)$\Rightarrow$(2). Work over a coordinate ball $U$ with a frame, so that $E|_U = U \times \mathbb{K}^r$ and $\nabla = d + \theta$ with $\theta = \sum_i\theta_i\,dx^i$, a matrix of $1$-forms. We look
for sections $s \colon U \to \mathbb{K}^r$ with $\partial_is + \theta_is = 0$ for all $i$. Their graphs are submanifolds of $U \times \mathbb{K}^r$ whose tangent spaces are spanned by the
vector fields
\[
X_i = \frac{\partial}{\partial x^i} - \sum_{j,k}(\theta_i)^k_jv^j\,\frac{\partial}{\partial v^k}, \qquad i = 1, \ldots, n,
\]
where $v^1, \ldots, v^r$ are the coordinates of $\mathbb{K}^r$ (real and imaginary parts separately if $\mathbb{K} = \CC$). These span a distribution of rank $n$. The coefficients of
$X_j$ in the $\partial/\partial v$ directions are linear in $v$, and a direct computation of the brackets gives
\[
[X_i, X_j] = -\sum_{k,l}\bigl(\partial_i\theta_j - \partial_j\theta_i + \theta_i\theta_j - \theta_j\theta_i\bigr)^k_lv^l\,\frac{\partial}{\partial v^k} = -\sum_{k,l}(F_{ij})^k_lv^l\,\frac{\partial}{\partial v^k},
\]
with $F_{ij} = F(\partial_i, \partial_j)$, where $F = \sum_{i<j}F_{ij}\,dx^i \wedge dx^j$: the terms $\partial_i$ applied to the coefficients of $X_j$ give $-\partial_i\theta_j$, and the
bracket of the two vertical parts gives the commutator of the matrices. So the distribution is involutive if and only if $F = 0$. If $F = 0$, the Frobenius theorem
(Theorem~\ref{smooth-manifolds:theorem-frobenius}) gives integral manifolds; each is transverse to the fibres $\{x\} \times \mathbb{K}^r$, hence locally the graph of a map $s \colon U' \to \mathbb{K}^r$, and the condition
that the graph is an integral manifold is exactly $\partial_is + \theta_is = 0$, that is, $\nabla s = 0$. Taking the integral manifolds through $(p, v_1), \ldots, (p, v_r)$ for a basis $(v_i)$ gives $r$ parallel
sections near $p$, which form a frame because they are linearly independent at $p$, hence nearby.

(2)$\Rightarrow$(3). Let $H \colon [0,1]^2 \to M$ be a homotopy with fixed endpoints from $\gamma_0$ to $\gamma_1$. Cover the square by finitely many small squares each mapped by $H$ into a neighbourhood
carrying a flat frame, as in the proof of Theorem~\ref{complex-analysis:theorem-cauchy-homotopy}. In a flat frame parallel transport along any path is the identity matrix, so the transport along the two
edge paths of each small square agree; moving across the squares one at a time shows that the transports along $\gamma_0$ and $\gamma_1$ coincide.

(3)$\Rightarrow$(2). On a simply connected neighbourhood $U$ of a point $p$, define $s_i(q) = P_\gamma(v_i)$ for a path $\gamma$ from $p$ to $q$ in $U$ and a basis $(v_i)$ of $E_p$; this is well defined by (3),
smooth by smooth dependence of solutions of differential equations, and gives a frame of parallel sections, in which the connection form vanishes.

For the last statement, parallel sections form a sheaf that is locally constant by (2), hence a local system, and a local frame of parallel sections is a frame of $E$, which gives the isomorphism with
$\mathcal{E}$. Conversely a local system $\mathcal{L}$ determines the bundle $E = \mathcal{L} \otimes_{\RR_M} C^\infty_M$ with the connection $d$ in any flat frame; these constructions are inverse to each
other. The correspondence with representations of $\pi_1$ is Theorem~\ref{sheaves:theorem-local-systems-representations}, and the monodromy of the local system of parallel sections is given by parallel
transport; the inverse appears because monodromy was defined as a right action.
\end{proof}

In a flat frame the transition functions between two flat frames are locally constant (both are parallel), so a flat bundle is one whose cocycle can be chosen locally
constant; this is the situation of Remark~\ref{connections:remark-no-canonical-d}, in which componentwise differentiation is well defined.

\begin{example}\label{connections:example-flat-circle}
Flat connections on the trivial line bundle over $S^1 = \RR/2\pi\ZZ$: every connection is $d + a(\theta)\,d\theta$, it is flat, and parallel transport around the circle is multiplication by $\exp(-\int_0^{2\pi}a)$. Two such connections are
isomorphic (by a gauge transformation $g \colon S^1 \to \CC^\times$, which changes $a\,d\theta$ by $g^{-1}dg$) if and only if their holonomies agree, because the $1$-forms $g^{-1}dg$ are exactly those with $\int_{S^1}g^{-1}dg \in 2\pi i\ZZ$
(a nonvanishing function has a continuous logarithm along the circle whose increase is $2\pi i$ times its winding number). So flat line bundles with connection on $S^1$ are classified by the holonomy in $\CC^\times$, in accordance with
Theorem~\ref{connections:theorem-flat-local-systems}, since $\pi_1(S^1) = \ZZ$.
\end{example}

\begin{counterexample}\label{connections:counterexample-flat-holonomy}
A flat connection can have nontrivial holonomy, and a bundle with trivial holonomy around contractible loops need not be trivial.
\begin{enumerate}
\item In Example~\ref{connections:example-flat-circle}, $d + a\,d\theta$ with $\int_0^{2\pi}a \notin 2\pi i\ZZ$ is flat but has no global parallel section: the holonomy is not $1$. Flatness
only gives parallel frames \emph{locally}.
\item The M\"obius bundle over $S^1$ (Example~\ref{vector-bundles:example-bundles}) has a flat connection, as its transition function $-1$ is locally constant, with holonomy $-1$
around the circle; it is not trivial.
\item Curvature is a property of the connection, not of the bundle: the trivial line bundle on $\RR^2$ carries the flat connection $d$ and the connection $d + x\,dy$ of
curvature $dx \wedge dy$. Only the topology of the bundle constrains the curvature, through characteristic classes (Chapter~\ref{characteristic-classes}); for instance the
tautological bundle of $\CC\PP^1$ admits no flat connection (Exercise~\ref{connections:exercise-flat-not-trivial}).
\end{enumerate}
\end{counterexample}

\section{Metric connections}\label{connections:section-metric}

\begin{definition}\label{connections:definition-metric-connection}
Let $E$ carry a Euclidean metric $\langle -, - \rangle$ (or a Hermitian metric in the complex case). A connection is \emph{metric} if
$d\langle s, t \rangle = \langle \nabla s, t \rangle + \langle s, \nabla t \rangle$, that is, $X\langle s, t\rangle = \langle\nabla_Xs, t\rangle + \langle s, \nabla_Xt\rangle$ for all
vector fields $X$ and sections $s$, $t$.
\end{definition}

\begin{proposition}\label{connections:proposition-metric-connections}
Metric connections exist, and $\nabla$ is metric if and only if parallel transport is an isometry, if and only if in every orthonormal local frame the connection form is skew-symmetric (skew-Hermitian in
the complex case); then the curvature takes values in the skew-symmetric (skew-Hermitian) endomorphisms, that is, in the Lie algebra of the structure group.
\end{proposition}

\begin{proof}
Existence: patch the connections with connection form $0$ in orthonormal local frames with a partition of unity, as in Proposition~\ref{connections:proposition-connections-exist}(4); each of them is
metric, since in an orthonormal frame $\langle s, t\rangle = \sum_is^i\bar t^i$ and $d\langle s, t\rangle = \sum_i(ds^i\,\bar t^i + s^i\,d\bar t^i)$, and the metric condition is preserved by
combinations with coefficients adding up to $1$. If $\nabla$ is metric and $s$, $t$ are parallel along $\gamma$, then
$\frac{d}{dt}\langle s, t \rangle = \langle\nabla_{\dot\gamma}s, t\rangle + \langle s, \nabla_{\dot\gamma}t\rangle = 0$ (the metric condition along the curve follows from the one on $M$ by
pulling back), so transport is an isometry. Conversely, suppose transport is an isometry, and let $v \in T_pM$, $\gamma$ a curve with $\dot\gamma(0) = v$, and $(u_i)$ a parallel frame along $\gamma$; the
$\langle u_i, u_j\rangle$ are constant. Writing $s \circ \gamma = \sum_ia^iu_i$ and $t \circ \gamma = \sum_jb^ju_j$, we get $\nabla_{\dot\gamma}s = \sum_i\dot a^iu_i$, and
$\frac{d}{dt}\langle s, t\rangle = \sum_{i,j}(\dot a^i\bar b^j + a^i\dot{\bar b}^j)\langle u_i, u_j\rangle = \langle\nabla_{\dot\gamma}s, t\rangle + \langle s, \nabla_{\dot\gamma}t\rangle$, which at
$t = 0$ is the metric condition in the direction $v$. In an orthonormal frame, $\langle e_i, e_j\rangle = \delta_{ij}$, and the
metric condition for $s = e_j$, $t = e_i$ reads $0 = \theta^i_j + \overline{\theta^j_i}$. Finally, if $\theta$ is skew-symmetric (skew-Hermitian), so are $d\theta$ and
$\theta \wedge \theta$, as $(\theta \wedge \theta)^T = -\theta^T \wedge \theta^T$; hence so is $F = d\theta + \theta \wedge \theta$.
\end{proof}

\begin{example}\label{connections:example-metric-projection}
The projection connection $\nabla s = P(ds)$ of a subbundle $E \subseteq M \times \RR^N$ (Example~\ref{connections:example-connections}(4)) is metric for the metric induced from $\RR^N$:
for sections $s$, $t$ of $E$, $d\langle s, t\rangle = \langle ds, t\rangle + \langle s, dt\rangle = \langle P\,ds, t\rangle + \langle s, P\,dt\rangle$, because $t = Pt$, $s = Ps$ and $P$ is
self-adjoint. So parallel transport on the sphere is by rotations, as in Example~\ref{connections:example-transport}(2), and the connection form of
Example~\ref{connections:example-connection-forms}(1) is skew-symmetric. On the tangent bundle of a Riemannian manifold there are many metric connections; among them exactly one is
\emph{torsion-free}, $\nabla_XY - \nabla_YX = [X, Y]$, the Levi-Civita connection (Chapter~\ref{riemannian-geometry}).
\end{example}

\begin{example}\label{connections:example-chern-connection}
On a complex vector bundle with a Hermitian metric over a complex manifold, there is a unique connection that is metric and compatible with the holomorphic structure ($\nabla^{0,1} = \bar\partial$): the
\emph{Chern connection} (Chapter~\ref{holomorphic-bundles}). Its curvature is a $(1,1)$-form, and for line bundles $\frac{i}{2\pi}F$ represents the first Chern class in de Rham cohomology, which is the
differential-geometric form of Theorem~\ref{vector-bundles:theorem-line-bundles}. The projection connection on the tautological line bundle of $\CC\PP^1$
(Example~\ref{connections:example-curvature-computations}(3)) is of this kind.
\end{example}

\section{Principal bundles}\label{connections:section-principal}

Vector bundles are often better understood through their bundles of frames, on which the structure group acts. This point of view is used in Chern--Weil theory and in gauge
theory; we only sketch it.

\begin{definition}\label{connections:definition-principal-bundle}
Let $G$ be a Lie group. A \emph{principal $G$-bundle} is a smooth surjection $\pi \colon P \to M$ with a free right action of $G$ preserving the fibres and acting simply transitively on them, which is locally
trivial: every point of $M$ has a neighbourhood $U$ with a $G$-equivariant diffeomorphism $\pi^{-1}(U) \cong U \times G$ over $U$. As for vector bundles, principal bundles are classified by cocycles
$g_{ij} \colon U_i \cap U_j \to G$, that is, by $\check{H}^1(M, \underline{G})$.
\end{definition}

\begin{example}\label{connections:example-principal}
The frame bundle $\mathrm{Fr}(E)$ of a rank $r$ real vector bundle, whose fibre over $p$ is the set of bases of $E_p$, is a principal $\GL_r(\RR)$-bundle, and $E$ is recovered as the associated bundle
$\mathrm{Fr}(E) \times_{\GL_r}\RR^r$; a metric reduces the structure group to $O(r)$, an orientation to $\GL_r^+$, and a complex structure to $\GL_{r/2}(\CC)$. Covering spaces are principal bundles with
discrete structure group. The oriented orthonormal frames of $TS^2$ form a principal $\mathrm{SO}(2)$-bundle which is diffeomorphic to $\mathrm{SO}(3)$: an oriented orthonormal
frame $(v, w)$ at $x \in S^2$ is the same as the rotation matrix with columns $x, v, w$. The Hopf map $S^3 \to S^2 = \CC\PP^1$, $(z_0, z_1) \mapsto [z_0 : z_1]$, is a principal
$U(1)$-bundle, the unit circle bundle of the tautological line bundle.
\end{example}

\begin{remark}\label{connections:remark-principal-connections}
A connection on a principal bundle is a $G$-equivariant splitting of $0 \to \ker d\pi \to TP \to \pi^*TM \to 0$, equivalently a $\mathfrak{g}$-valued $1$-form $\omega$ on $P$ restricting to the
Maurer--Cartan form on fibres and satisfying $R_g^*\omega = \Ad(g^{-1})\omega$. Its curvature is $\Omega = d\omega + \frac12[\omega, \omega]$, and connections on $P$ correspond to connections on all
associated vector bundles, compatible with the structure group. This is the framework of gauge theory and of Chern--Weil theory, which we use in Chapter~\ref{characteristic-classes}; see
\cite[Chapters I--II]{kobayashi-nomizu}. For a vector bundle, a local frame is a local section of the frame bundle, and pulling back $\omega$ by it gives the connection form
$\theta$ of Proposition~\ref{connections:proposition-connections-exist}; the transformation rule $\theta^{\prime} = g^{-1}\theta g + g^{-1}dg$ is the pullback of the equivariance of
$\omega$.
\end{remark}

\section{Exercises}\label{connections:section-exercises}

\begin{exercise}\label{connections:exercise-curvature-sphere}
Compute the connection and curvature forms of the Levi-Civita connection of the sphere of radius $\rho$ in $\RR^3$, that is, of the projection connection of its tangent
bundle, in the frame of Example~\ref{connections:example-connection-forms}(1), and verify that the curvature is $F = \rho^{-2}\,dA \otimes J$, where $dA$ is the area form, so that
the integral of the Gaussian curvature $\rho^{-2}$ is $4\pi$ (compare the Gauss--Bonnet theorem, Chapter~\ref{characteristic-classes}).
\end{exercise}

\begin{solution}
Scaling by $\rho$ does not change the unit vectors $e_1$, $e_2$ or the normal, so the computation of Example~\ref{connections:example-connection-forms}(1) is unchanged:
$\theta^2_1 = \cos\vartheta\,d\varphi$ and $F^1_2 = \sin\vartheta\,d\vartheta \wedge d\varphi$. The area form is now $dA = \rho^2\sin\vartheta\,d\vartheta \wedge d\varphi$, so
$F = \rho^{-2}dA \otimes J$, and $\int\rho^{-2}\,dA = \rho^{-2} \cdot 4\pi\rho^2 = 4\pi$.
\end{solution}

\begin{exercise}\label{connections:exercise-flat-not-trivial}
Give an example of a flat bundle that is not trivial, and of a bundle admitting no flat connection.
\end{exercise}

\begin{solution}
The M\"obius bundle over $S^1$ is flat (its transition function is locally constant) but not trivial
(Exercise~\ref{vector-bundles:exercise-mobius}). The tautological bundle $\mathcal{O}(-1)$ on $\CC\PP^1$ admits no flat connection: by Theorem~\ref{connections:theorem-flat-local-systems} a flat bundle on a
simply connected manifold is determined by a representation of the trivial group, hence is trivial, while $\mathcal{O}(-1)$ is not trivial because $c_1(\mathcal{O}(-1)) \neq 0$
(Example~\ref{vector-bundles:example-line-bundles}).
\end{solution}

\begin{exercise}\label{connections:exercise-holonomy-abelian}
Show that for a line bundle with connection, the holonomy is given by $\exp(-\int_\Sigma F)$ for a loop bounding a surface $\Sigma$, and deduce that flat line bundles on a simply connected manifold are
trivial.
\end{exercise}

\begin{solution}
For a disc this is Proposition~\ref{connections:proposition-line-holonomy}; for a surface $\Sigma$ with one boundary circle over which the pulled back bundle is trivial the same
proof applies. If $F = 0$ and $M$ is simply connected, the holonomy is trivial, since every loop bounds a disc (up to homotopy, which does not change the transport of a flat
connection by Theorem~\ref{connections:theorem-flat-local-systems}); so transport from a fixed point along any path gives a well-defined nowhere vanishing parallel section, which
trivialises the bundle.
\end{solution}

\begin{exercise}\label{connections:exercise-induced-connections}
Verify that a connection on $E$ induces connections on $E^\vee$ and $E \otimes F$, and that the curvature of $E \otimes F$ is $F_E \otimes 1 + 1 \otimes F_F$.
\end{exercise}

\begin{solution}
This is Proposition~\ref{connections:proposition-induced-connections}(1),(3) and Proposition~\ref{connections:proposition-pullback-tensor}(2). With vector fields, for
$s \otimes s^{\prime}$: $\nabla_X\nabla_Y(s \otimes s^{\prime}) = \nabla_X\nabla_Ys \otimes s^{\prime} + \nabla_Ys \otimes \nabla_Xs^{\prime} + \nabla_Xs \otimes \nabla_Ys^{\prime} + s \otimes \nabla_X\nabla_Ys^{\prime}$; the
middle terms are symmetric in $X$ and $Y$ and cancel in $F(X, Y)$, which leaves $F_E(X, Y)s \otimes s^{\prime} + s \otimes F_F(X, Y)s^{\prime}$.
\end{solution}

\begin{exercise}\label{connections:exercise-projection-curvature}
Let $E \subseteq M \times \RR^N$ be a subbundle with orthogonal projection $P$ and the projection connection $\nabla s = P\,ds$. Show that $P\,dP\,P = 0$ and that the curvature is
$F(s) = P\,dP \wedge dP\,s$ for sections $s$ of $E$, where $dP$ is the matrix of $1$-forms of the derivatives of $P$. Recover Example~\ref{connections:example-curvature-computations}(2).
\end{exercise}

\begin{solution}
Differentiating $P^2 = P$ gives $dP = dP\,P + P\,dP$; multiplying by $P$ on both sides gives $P\,dP\,P = 2P\,dP\,P$, so $P\,dP\,P = 0$. For a section $s = Ps$, extend $\nabla$ to
$E$-valued $1$-forms by $d_\nabla\omega = P\,d\omega$; then $F(s) = P\,d(P\,ds) = P\,dP \wedge ds$, and $ds = d(Ps) = dP\,s + P\,ds$, so
$F(s) = P\,dP \wedge dP\,s + P\,dP\,P \wedge ds = P\,dP \wedge dP\,s$. For the sphere, $P = I - xx^T$ and $dP = -dx\,x^T - x\,dx^T$, and evaluating on $e_1$, $e_2$ gives the
formula of the example.
\end{solution}

\begin{exercise}\label{connections:exercise-gauge-plane}
On the trivial real line bundle over $\RR^2$, show that the connections $d + x\,dy$ and $d - y\,dx$ have the same curvature, and find a function $g \colon \RR^2 \to \RR_{>0}$
such that the frame $e^{\prime} = ge$ transforms one connection form into the other.
\end{exercise}

\begin{solution}
Both curvatures are $d(x\,dy) = dx \wedge dy = d(-y\,dx)$. In the frame $e^{\prime} = ge$ the connection form $x\,dy$ becomes $x\,dy + g^{-1}dg = x\,dy + d\log g$
(Proposition~\ref{connections:proposition-connections-exist}(3), with commuting $1 \times 1$ matrices). We need $d\log g = -y\,dx - x\,dy = -d(xy)$, so $g = e^{-xy}$.
\end{solution}

\begin{exercise}\label{connections:exercise-torsion}
For a connection $\nabla$ on $TM$ define the \emph{torsion} $T(X, Y) = \nabla_XY - \nabla_YX - [X, Y]$. Show that $T$ is a tensor, a $TM$-valued $2$-form, and that in coordinates,
with $\nabla_{\partial_i}\partial_j = \sum_k\Gamma^k_{ij}\partial_k$, it vanishes if and only if $\Gamma^k_{ij} = \Gamma^k_{ji}$. Show that the projection connection on $TS^n$ is
torsion-free.
\end{exercise}

\begin{solution}
$T(fX, Y) = f\nabla_XY - (Yf)X - f\nabla_YX - f[X, Y] + (Yf)X = fT(X, Y)$, and $T$ is antisymmetric, so it is $C^\infty$-linear in both arguments and is a tensor by
Lemma~\ref{connections:lemma-tensoriality}. On coordinate fields $[\partial_i, \partial_j] = 0$ and $T(\partial_i, \partial_j) = \sum_k(\Gamma^k_{ij} - \Gamma^k_{ji})\partial_k$. For $TS^n$,
$\nabla_XY - \nabla_YX = P(X(Y) - Y(X)) = P([X, Y]) = [X, Y]$, where $X(Y)$ is the derivative in $\RR^{n+1}$ of the extended fields and the bracket of fields tangent to $S^n$ is
tangent to it.
\end{solution}

\begin{exercise}\label{connections:exercise-connections-on-intervals}
Show that every vector bundle with connection over an interval $I \subseteq \RR$ has a frame of parallel sections, and deduce again that every vector bundle over an interval
is trivial.
\end{exercise}

\begin{solution}
Fix $t_0 \in I$ and a basis $(v_i)$ of the fibre over $t_0$. The parallel sections along the curve $t \mapsto t$ with initial values $v_i$ exist on all of $I$
(Proposition~\ref{connections:proposition-parallel-transport}, applied on compact subintervals) and are linearly independent everywhere, since parallel transport is an isomorphism;
so they form a global frame. The curvature plays no role: on a $1$-dimensional base every connection is flat.
\end{solution}
